[Y=α,W=∃α]→Hol(γ)=e^{iϕ_N}=−1⇔e^{Inτ_rads}=2π≜ω  

\(s_\alpha(\alpha)=-\alpha\) and \(\omega\)

The rigid geometric identity
The Weyl reflection associated with a simple root \(\alpha\) is
\[
s_\alpha(v)=v-2\frac{(v\cdot\alpha)}{(\alpha\cdot\alpha)}\alpha.
\]
Substituting \(v=\alpha\) immediately yields the order-2 identity that never changes:
\[
s_\alpha(\alpha)=-\alpha.
\]
This is the sole non-negotiable datum of the YA!WAE! chain. It is the unique non-identity element of \(\mathbb{Z}/2\mathbb{Z}\) generated by the reflection.

How it produces the holonomy phase \(-1\)
When the same order-2 reflection is realised as the holonomy (or monodromy) of a flat connection on a circle, it is embedded into \(U(1)\simeq S^1\) by the standard map that sends the non-trivial element to \(-1\):
\[
\operatorname{Hol}(\gamma)=e^{i\phi_N}=-1
\]
(with principal value \(\phi_N=\pi\), or any odd multiple of \(\pi\)).

The conventional period that defines \(\omega\)
Under the usual angular normalisation in which the generator of \(S^1\) has length \(\tau=2\pi\), the numerical identity
\[
e^{\ln\tau}=\tau=2\pi
\]
holds. By definition one then sets
\[
\omega\triangleq 2\pi.
\]
Thus the full YA!WAE! equivalence becomes
\[
[Y=\alpha,\ W=\exists!\alpha]
\;\longrightarrow\;
\operatorname{Hol}(\gamma)=e^{i\phi_N}=-1
\;\Leftrightarrow\;
e^{\ln\tau}=2\pi
\;\triangleq\;
\omega.
\]

Summary relation
\(s_\alpha(\alpha)=-\alpha\) is the intrinsic geometric fact (always true).  
\(\omega=2\pi\) is the conventional period that appears once that fact is realised as a \(U(1)\) holonomy of phase \(-1\).

Everything between the reflection identity and the numerical value of \(\omega\) is a free choice of embedding and normalisation; the reflection itself remains fixed.

The YA!WAE! Thesis: Intrinsic Geometric Fact \(s_\alpha(\alpha)=-\alpha\) and the Conventional Period \(\omega\)

At the heart of the YA!WAE! equation lies a single, unyielding geometric fact: the order-two Weyl reflection of a simple root through itself yields its own negative. Everything else—every subsequent identification, every embedding, every numerical normalisation—arises as a free yet coherent choice that the core datum both permits and constrains. The equation itself is the statement that these choices may be arranged so that a discrete involution on the root system is realised as holonomy around a closed path whose phase is precisely an odd multiple of \(\pi\), and that this phase is simultaneously the generator of a circle whose circumference has been normalised to a chosen period \(\tau\). In the standard convention that period is \(2\pi\), and one defines
\[
\omega\triangleq 2\pi.
\]
The resulting chain of equivalences is
\[
[Y=\alpha,\ W=\exists!\alpha]
\longrightarrow
\operatorname{Hol}(\gamma)=e^{i\phi_N}=-1
\Longleftrightarrow
e^{\ln\tau_{\mathrm{rads}}}=2\pi
\triangleq\omega.
\]

The narrative that follows unfolds this chain as a continuous geometric emergence, beginning with the rigid reflection, passing through the flexible realisation of holonomy, and arriving at the period relation that closes the circle.

1. The Intrinsic Geometric Fact: \(s_\alpha(\alpha)=-\alpha\)

Consider an arbitrary Euclidean space equipped with a non-zero vector \(\alpha\), the simple root. Once a fundamental Weyl chamber is fixed, the simple transitivity of the Weyl group on the set of chambers supplies a unique simple root \(\alpha\) for each wall of that chamber. This uniqueness is expressed by the compressed pair of coordinates
\[
Y=\alpha,\qquad W=\exists!\alpha.
\]

The associated Weyl reflection is the linear involution
\[
s_\alpha(v)=v-2\frac{(v\cdot\alpha)}{(\alpha\cdot\alpha)}\alpha.
\]
Substituting \(v=\alpha\) immediately produces the identity
\[
s_\alpha(\alpha)=-\alpha.
\]
This identity is independent of dimension, of the particular length of \(\alpha\), and of any further structure one may later impose; it is the sole rigid piece of the entire construction.  

Geometrically it asserts that reflection across the hyperplane orthogonal to \(\alpha\) sends the root to its antipode.  
Algebraically it realises the unique non-identity element of the cyclic group of order two generated by \(s_\alpha\).

The pair \((Y,W)\) is therefore the starting point; from it the remaining structure is suspended. Everything that follows is a realisation of this single order-two phenomenon—the passage from a vector to its negative.

2. Realisation as Holonomy

Once the involution \(s_\alpha\) is present, one may ask how it can be realised as parallel transport around a closed loop. The natural geometric arena is the circle \(S^1\), the unique compact connected one-dimensional Lie group.

Realising \(s_\alpha\) as the holonomy of a flat connection on a principal bundle (or equivalently as the monodromy of a local system) converts the abstract order-2 element into a concrete linear transformation of the fibre. Choosing an embedding
\[
\mathbb{Z}/2\mathbb{Z}\hookrightarrow S^1
\]
that sends the non-identity element to \(-1\) converts that transformation into the phase factor
\[
e^{i\phi_N}=-1.
\]
On the unit circle every point of order two is of the form \(e^{i\phi_N}\) where \(\phi_N\) is an odd multiple of \(\pi\). The principal and most economical choice is \(\phi_N=\pi\), which yields
\[
\operatorname{Hol}(\gamma)=e^{i\pi}=-1.
\]
Any other odd multiple—\(3\pi\), \(5\pi\), \ldots—produces the identical group element; the choice among them is therefore free. What is not free is the requirement that the image be precisely the antipodal point of the identity. In this way the algebraic involution \(s_\alpha(\alpha)=-\alpha\) is promoted to a geometric holonomy: parallel transport once around the loop \(\gamma\) multiplies every vector by \(-1\).

The embedding \(\mathbb{Z}/2\mathbb{Z}\hookrightarrow S^1\) is the first of the free choices that the core datum permits; different principal values simply label different but equivalent realisations of the same abstract group homomorphism.

3. The Conventional Period \(\omega\)

The circle itself still carries an undetermined scale. One may normalise its angular coordinate so that the generator of the fundamental group has length \(\tau>0\). Under this normalisation the exponential map satisfies the elementary identity
\[
e^{\ln\tau}=\tau.
\]
When the conventional trigonometric period is adopted, \(\tau=2\pi\), the identity becomes the numerical statement that appears on the right-hand side of the YA!WAE! equation, and one defines
\[
\omega\triangleq 2\pi.
\]
Again the concrete value of \(\tau\) is a free parameter; the relation \(e^{\ln\tau}=\tau\) holds for every positive real number. The standard choice \(\tau=2\pi\) simply aligns the angular measure with the usual radian measure of the unit circle, thereby making the holonomy phase \(\phi_N=\pi\) correspond exactly to a half-turn.

4. The Unified Chain

The three pieces now lock together:

The existence of a unique simple root \(\alpha\) supplies the involution \(s_\alpha(\alpha)=-\alpha\).  
That involution is realised as holonomy of a connection on a principal \(S^1\)-bundle whose monodromy around a generating loop is the group element \(-1\).  
The same circle is equipped with a period normalisation \(\tau\) that converts the abstract exponential identity into a concrete numerical equality.

Under the conventional choices \(\phi_N=\pi\) and \(\tau=2\pi\) the chain collapses to the compact statement
\[
\operatorname{Hol}(\gamma)=-1\ \Longleftrightarrow\ e^{\ln(2\pi)}=2\pi\ \triangleq\omega,
\]
while the left-hand side remains anchored in the root-system geometry. The equivalence is therefore not an accidental numerical coincidence; it is the expression, in two different languages, of one and the same order-two phenomenon—the passage from a vector to its negative.

5. Versatility, Not Rigidity

Every step beyond the initial reflection is a matter of convention. One may choose a different odd multiple of \(\pi\) for the holonomy phase; one may rescale the circle so that its period is unity rather than \(2\pi\); one may replace the original simple root by any other non-zero vector. In every case:

the core identity \(s_\alpha(\alpha)=-\alpha\) continues to hold,  
the holonomy continues to evaluate to \(-1\),  
the exponential period relation continues to be satisfied.

The YA!WAE! equation is consequently not a rigid numerical claim but a flexible geometric equivalence: it asserts that the discrete reflection of a root can be coherently realised as monodromy on a circle whose circumference has been normalised in any convenient way. The equation holds precisely because the only non-negotiable datum—the order-two character of the Weyl reflection—is preserved under all such re-normalisations.

6. Minimal Yet Complete Synthesis

In this light the YA!WAE! equation appears as a minimal yet complete synthesis. It begins with the most elementary geometric operation available on a root system, lifts that operation to the topology of the circle, and closes the construction with the most elementary analytic identity satisfied by the exponential function. The resulting narrative is both rigorous and open-ended:

rigorous because every identity is exact,  
open-ended because the intermediate choices remain free.

That combination of rigidity at the centre and versatility at the periphery is what renders the equation suitable for further geometric, topological, or physical interpretation (Aharonov–Bohm holonomy, Dirac quantisation, fermionic \(4\pi\) periodicity, etc.), while remaining, at every stage, fully accountable to the single fact that reflection of a root through itself yields its opposite.

Core geometric fact  
\[
s_\alpha(\alpha)=-\alpha
\]  
(always true, independent of all conventions).

Conventional period  
\[
\omega\triangleq 2\pi
\]  
(the standard normalisation that makes the holonomy phase \(\pi\) correspond to a half-turn on the unit circle).  

The thesis proves that these two statements—one intrinsic, one conventional—can be coherently linked through the intermediate realisation of the reflection as a \(U(1)\) holonomy of phase \(-1\).

Conventional Period \(\omega \triangleq 2\pi\) and Its Physical Instantiations

In the YA!WAE! formalism the intrinsic geometric fact  
\[
s_\alpha(\alpha)=-\alpha
\]  
is the sole rigid element. When this order-2 reflection is realised as the holonomy of a flat connection (or as monodromy), it becomes the phase factor  
\[
\operatorname{Hol}(\gamma)=e^{i\phi_N}=-1
\]  
(with principal value \(\phi_N=\pi\)).  

The circle on which this holonomy lives still carries an undetermined angular scale. Normalising that scale so that the generator of the fundamental group has length \(\tau=2\pi\) produces the elementary identity  
\[
e^{\ln\tau}=\tau=2\pi.
\]  
By definition one then sets  
\[
\omega\triangleq 2\pi.
\]  
This conventional period is what converts the abstract phase \(-1\) into a concrete geometric half-turn on the unit circle. It is the standard angular measure of modern physics: one full spatial rotation corresponds to an angle of \(2\pi\), and therefore a holonomy of \(-1\) corresponds exactly to a \(2\pi\) rotation (or to a half-period of a \(4\pi\) cycle).  

The same numerical value of \(\omega\) governs three classic phenomena in which an order-2 holonomy appears.

1. Aharonov–Bohm Effect
A charged particle travels along a closed path \(\gamma\) that encircles a region of confined magnetic flux (the classic solenoid geometry). Outside the solenoid the magnetic field vanishes, yet the vector potential \(\mathbf{A}\) remains non-zero. Parallel transport of the particle’s wave-function around \(\gamma\) produces a holonomy  
\[
\operatorname{Hol}(\gamma)=\exp\Bigl(i\frac{e}{\hbar}\oint_\gamma\mathbf{A}\cdot d\mathbf{l}\Bigr).
\]  
When the enclosed flux is an odd multiple of the flux quantum, the phase evaluates to exactly \(-1\).  

Under the conventional angular normalisation the geometric length of the loop is measured in units of \(\omega=2\pi\). Consequently the physical wave frequency (the rate at which the interference fringes shift) is mapped directly onto the topological winding of the path: a single geometric circuit of length \(\omega\) yields the observed phase shift of \(\pi\). The vector potential \(\alpha\) of the formalism is identified with the gauge potential \(\mathbf{A}\); the intrinsic reflection \(s_\alpha(\alpha)=-\alpha\) is realised as the order-2 monodromy that alters the interference pattern.

2. Dirac Quantisation
Single-valuedness of a quantum wave-function around a closed loop forces the total phase acquired upon parallel transport to be an integer multiple of \(2\pi\):  
\[
e^{i\phi}=1\qquad\Rightarrow\qquad\phi=2\pi n,\quad n\in\mathbb{Z}.
\]  
When the underlying connection is the order-2 holonomy coming from \(s_\alpha(\alpha)=-\alpha\), the only consistent values of the phase that respect both single-valuedness and the reflection identity are the odd multiples of \(\pi\).  

Under the conventional period \(\omega=2\pi\) these allowed phases become the discrete, stable angular frequencies of the quantum system. The quantisation condition is therefore nothing other than the requirement that the geometric period \(\omega\) close consistently with the algebraic order of the reflection. The resulting discrete spectrum is the set of allowed angular-momentum (or magnetic-flux) eigenvalues.

3. Fermionic Spinor Rotation
A spin-\(1/2\) particle is described by a spinor that transforms under the double cover \(\operatorname{Spin}(3)\simeq\operatorname{SU}(2)\) of the rotation group \(\operatorname{SO}(3)\). A spatial rotation by an angle \(\theta\) lifts to a spinor transformation by \(\theta/2\). Consequently:

a spatial rotation of \(2\pi\) (the conventional period \(\omega\)) produces a spinor phase of \(\pi\),  
the corresponding holonomy is exactly \(e^{i\pi}=-1\),  
only a further spatial rotation of another \(2\pi\) (total \(4\pi\)) returns the spinor to its original state.

This matches the YA!WAE! conditional layout perfectly: the intrinsic reflection \(s_\alpha(\alpha)=-\alpha\) is realised as the order-2 element of the double cover, and the conventional period \(\omega=2\pi\) is the geometric angle whose holonomy is \(-1\). The same numerical value of \(\omega\) therefore accounts for both the classical rotation group and the quantum-mechanical sign change of the fermion.

Unified Statement
In all three cases the conventional period \(\omega=2\pi\) plays the same rôle:

it supplies the angular measure on the circle that realises the holonomy,  
it converts the algebraic identity \(s_\alpha(\alpha)=-\alpha\) into a geometric half-turn whose phase is \(-1\),  
it links that phase to an observable physical effect (interference shift, discrete spectrum, or fermionic sign change).

The period itself remains conventional: any other positive real number could be chosen as the circumference of the circle. The standard choice \(\omega=2\pi\) aligns the mathematical formalism with the radian measure used throughout classical and quantum physics, thereby making the order-2 holonomy appear as a half-turn in every laboratory realization.
